\documentclass{article}
\usepackage{graphicx}
\usepackage{amsmath}
\usepackage{amsfonts}
\usepackage{bookmark}
\usepackage{booktabs}
\usepackage{exercises}
\graphicspath{ {} }
\usepackage{indentfirst}
\usepackage{amsthm}
\usepackage{yhmath}
\usepackage{geometry}
\usepackage{placeins}
\usepackage{csquotes}
\usepackage{pgfplots}
\newcommand*{\oldneg}{\mathord{\sim}}
\newcommand\showdiv[1]{\overline{\smash{)}#1}}
\newcommand\Mydiv[1]{%
\smash{\raise.3ex\hbox{$\big)$}}\mkern-8mu
        \overline{\enspace\strut#1}}
\setlength{\tabcolsep}{5pt}
\theoremstyle{definition}
\newtheorem{definition}{Definition}

\setlength{\parindent}{.5in}
\setlength{\parskip}{12pt}
\renewcommand{\arraystretch}{2}
\geometry{margin=1in}
\newcommand{\?}{\stackrel{?}{=}}
\pgfplotsset{width=7cm,compat=1.18}\usepgfplotslibrary{statistics}

\title{Problem Solving Review (pg. 422)}
\begin{document}
\maketitle

\section*{1}
(a): The mean of the data set is 1.895.

(b): I did not know how to find the standard deviation.

\section*{2}
(a): The margin of error is $\pm\frac{1}{\sqrt{1022}}\approx 0.03$


(b): An interval that is likely to contain the exact percent of all online shoppers who shop because of the convenience is 
$70\%-76\%$.

\section*{3}
We can use the normal distribution model to find that the percent of trees at least 220 centimeters in height is $13.5+2.35+0.135=15.9855\%\approx 16\%$. 

\section*{4}
It is not really possible to do this with an observational study. It seems unlikely that you would be able to just find orchids with and without the high levels of 
dissolved calcium, so it is likely that the research needs to be conducted by applying treatment in an experiment. 
This can be done by selecting a randomized control group that does not 
get water with high levels of dissolved calcium, and a randomized treatment group which does. Then, observe the differences in height in the two groups at regular intervals.

\section*{5}
By the normal distribution curve, 
68\% of the amounts dispensed lie in the range $\bar{x}-\sigma$ to $\bar{x}+\sigma$, where $\bar{x}$ is the mean and $\sigma$ is the standard deviation. Thus, 
$68\%$ fall withing $10.5-0.75=9.75$ and $10.5+0.75=11.25$ ounces.

\section*{6}
Using the normal distribution, and the fact that the mean is 45 minutes and the standard deviation is 12 minutes, we can tell that the amount of shoppers spending 
between 45 and 69 minutes is in the range $\bar{x}$ to $\bar{x}+2\sigma$. Thus the probability is $34+13.5=47.5\%$.

\section*{7}
(a): 315 was 15\% of the respondents, meaning $\frac{315}{\frac{15}{100}}=2100$ students were surveyed.


(b): The margin of error is $\pm\frac{1}{\sqrt{2100}}\approx 0.02$.


(c): Thus, the interval that is likely to contain the exact percent is $13\%-17\%$.

\section*{8}
This is a convenience survey, as the station is just going to a nearby (easy to access) sports stadium. This is obvsiouly a biased sample, as people attending a sports 
game are more likely to watch sports on TV.

\end{document}
